% ============================================================
% ECEN 662 Project Report
% Compile: pdflatex proposal && bibtex proposal && pdflatex proposal && pdflatex proposal
% Requires neurips_2025.sty in the same directory
% ============================================================

\documentclass{article}
\usepackage[preprint]{neurips_2025}

\usepackage[utf8]{inputenc}
\usepackage[T1]{fontenc}
\usepackage{amsmath, amssymb, amsfonts}
\usepackage{algorithm}
\usepackage{algpseudocode}
\usepackage{booktabs}
\usepackage{graphicx}
\usepackage{natbib}
\usepackage{hyperref}

\newcommand{\E}{\mathbb{E}}
\newcommand{\KL}{\mathrm{KL}}

\title{The Many Faces of EM: Classical, Variational, Monte Carlo,
       and Online Expectation-Maximization on Hidden Markov Models}

\author{
  Brody Jordan \\
  Department of Electrical \& Computer Engineering \\
  Texas A\&M University \\
  \texttt{blj@tamu.edu}
}

\begin{document}
\maketitle

\begin{abstract}
  The Expectation-Maximization (EM) algorithm resolves the intractability of
  maximum likelihood estimation in latent variable models via successive
  approximation. When the E-step itself is intractable, three variants address
  the bottleneck differently: Variational EM approximates the posterior with a
  mean-field family, Monte Carlo EM estimates it via a particle filter, and
  Online EM processes data one sequence at a time. This report compares all four
  on Hidden Markov Models, where the exact Baum-Welch solution provides a
  verifiable ground truth.
\end{abstract}

% ============================================================
\section{Introduction}
% ============================================================

The central challenge in latent variable models is that the marginal
likelihood requires integrating out the unobserved variable $U$:
\begin{equation}
  f_\theta(x) = \int f_\theta(x, u)\, du.
  \label{eq:marginal}
\end{equation}
This integral appears \emph{inside} the log in the maximum likelihood estimation (MLE) objective, making
direct maximization intractable \citep{liu_ecen662_em}. EM resolves this via
successive approximation, but the E-step can itself become intractable for
complex models. This project studies four strategies for handling that
bottleneck on a shared Hidden Markov Model (HMM) testbed.

% ============================================================
\section{Background}
% ============================================================

\subsection{Maximum Likelihood Estimation with Latent Variables}

Given observations $X$ and latent variables $U$, the MLE is
\citep{liu_ecen662_em}:
\begin{equation}
  \hat\theta(x) \in \arg\max_{\theta \in \Theta}\; \log f_\theta(x),
  \qquad f_\theta(x) = \int f_\theta(x, u)\, du.
\end{equation}
The marginalization inside the log is the core difficulty.

\subsection{The Evidence Lower Bound (ELBO)}

For any auxiliary distribution $q(u)$, the log-likelihood decomposes as
\citep{liu_ecen662_em}:
\begin{equation}
  \log f_\theta(x)
  \;=\;
  \underbrace{
    \E_q\!\left[\log f_\theta(x, u)\right] - \E_q\!\left[\log q(u)\right]
  }_{\displaystyle\mathcal{L}(\theta,\, q)\ (\text{ELBO})}
  \;+\;
  \KL\!\left(q \;\|\; f_\theta(u \mid x)\right)
  \;\geq\; \mathcal{L}(\theta,\, q).
  \label{eq:elbo}
\end{equation}
Since $\KL \geq 0$, the ELBO is always a lower bound on the log-likelihood.
Setting $q = f_\theta(u \mid x)$ collapses the KL term to zero and makes the
bound tight.

\subsection{The Computational Trap}

For a Gaussian Mixture Model with $K$ components and $n$ observations, the
evidence requires summing over all $K^n$ discrete assignment configurations
\citep{liu_ecen662_vi}:
\begin{equation}
  m(x) = \sum_c p(c) \int f(\mu) \prod_{i=1}^n f(x_i \mid \mu_{c_i})\, d\mu.
\end{equation}
This scales exponentially in $n$, making exact inference strictly intractable.

% ============================================================
\section{Methods}
% ============================================================

All four methods alternate between an E-step (computing or approximating
the posterior over latent states) and an M-step (updating $\theta$). They
differ only in how the E-step is handled.

\subsection{Base EM}
\subsection{Variational EM (VEM)}
\subsection{Expectation conditional maximization (ECM)}

ECM \citep{meng1993maximum} addresses cases where the joint M-step maximization
over all parameters is hard, but maximizing over subsets conditional on the
others is tractable. Partition $\theta$ into subvectors
$(\theta_1, \ldots, \theta_p)$. Each iteration proceeds as:

\begin{algorithm}[ht]
  \caption{Expectation Conditional Maximization (ECM) \citep{meng1993maximum}}
  \begin{algorithmic}[1]
    \State Initialize $\theta^{(0)}$
    \Repeat
    \State \textbf{E-step:} Compute
    $Q(\theta,\, \theta^{(t)}) =
      \mathbb{E}_{f(u \mid x,\, \theta^{(t)})}\!\left[\log f_\theta(x, u)\right]$
    \For{$j = 1, \ldots, p$}
    \State \textbf{CM-step $j$:}
    $\theta_j^{(t+1)} = \arg\max_{\theta_j}\;
      Q\!\left(\theta_j,\, \theta_{-j}^{(t+1)},\, \theta^{(t)}\right)$
    \EndFor
    \Until{$\left|\log f_{\theta^{(t+1)}}(x) - \log f_{\theta^{(t)}}(x)\right| < \varepsilon$}
  \end{algorithmic}
\end{algorithm}


\subsection{Generalized EM (GEM)}

% \subsection{The EM Algorithm}

% The E-step computes the exact posterior marginals via the forward-backward
% algorithm \citep{liu_ecen662_em}. Define:
% \begin{align}
%   \alpha_t(k) & = p(x_1, \ldots, x_t,\, Z_t = k),       \\
%   \beta_t(k)  & = p(x_{t+1}, \ldots, x_T \mid Z_t = k).
% \end{align}
% The posterior state marginals and pairwise marginals are:
% \begin{align}
%   \gamma_t(k) & = \frac{\alpha_t(k)\,\beta_t(k)}{\sum_j \alpha_t(j)\,\beta_t(j)},
%   \label{eq:gamma}                                                                      \\
%   \xi_t(j, k) & \propto \alpha_t(j)\,A_{jk}\,p(x_{t+1} \mid Z_{t+1}=k)\,\beta_{t+1}(k).
%   \label{eq:xi}
% \end{align}
% The M-step updates $\theta$ in closed form using $\gamma_t(k)$ and $\xi_t(j,k)$:
% \begin{equation}
%   \hat\mu_k = \frac{\sum_t \gamma_t(k)\,x_t}{\sum_t \gamma_t(k)}, \qquad
%   \hat A_{jk} = \frac{\sum_t \xi_t(j,k)}{\sum_t \gamma_t(j)}.
%   \label{eq:mstep}
% \end{equation}

% \begin{algorithm}[ht]
%   \caption{The EM Algorithm}
%   \begin{algorithmic}[1]
%     \State Initialize $\theta^{(0)}$
%     \Repeat
%     \State \textbf{E-step:} Run forward-backward to compute
%     $\gamma_t(k)$, $\xi_t(j,k)$ for all $t$, $k$
%     \State \textbf{M-step:} Update $\theta$ via Eq.~\eqref{eq:mstep}
%     \Until{$|\log f_{\theta^{(t+1)}}(x) - \log f_{\theta^{(t)}}(x)| < \varepsilon$}
%   \end{algorithmic}
% \end{algorithm}

% \subsection{Variational EM}

% When the exact posterior is intractable, VI replaces it with the closest
% member of a tractable family $\mathcal{Q}$ \citep{liu_ecen662_vi}:
% \begin{equation}
%   q^* = \arg\min_{q \in \mathcal{Q}}\; \KL\!\left(q(z_{1:T}) \;\|\; f_\theta(z_{1:T} \mid x)\right).
%   \label{eq:vi}
% \end{equation}
% We use the mean-field family $\mathcal{Q} = \{q : q(z_{1:T}) = \prod_t q_t(z_t)\}$.
% The coordinate-ascent update for each factor is \citep{liu_ecen662_vi}:
% \begin{equation}
%   \log q_t^*(z_t = k) \;\propto\;
%   \log p(x_t \mid k)
%   + \sum_j q_{t-1}(j) \log A_{jk}
%   + \sum_j q_{t+1}(j) \log A_{kj}.
%   \label{eq:mf}
% \end{equation}
% The M-step substitutes $q_t(k)$ for $\gamma_t(k)$ in Eq.~\eqref{eq:mstep}.

% \subsection{Monte Carlo EM (MCEM)}

% The E-step expectation is estimated with $M$ particles via a bootstrap
% particle filter. At each time step:
% \begin{align}
%   z_t^{(m)} & \sim \mathrm{Cat}\!\left(A_{z_{t-1}^{(m)},\cdot}\right), \\
%   w_t^{(m)} & \propto p\!\left(x_t \mid z_t^{(m)}\right), \quad
%   \hat\gamma_t(k) = \sum_{m=1}^M w_t^{(m)}\,\mathbf{1}\!\left[z_t^{(m)} = k\right].
%   \label{eq:mcem}
% \end{align}
% As $M \to \infty$, $\hat\gamma_t(k) \to \gamma_t(k)$ (the exact Baum-Welch
% marginal). The M-step is identical to Eq.~\eqref{eq:mstep}.

% \subsection{Online EM}

% Because the HMM complete-data log-likelihood belongs to the exponential
% family, the M-step only requires sufficient statistics $s(x, z)$.
% Online EM maintains a running average updated after each sequence
% \citep{liu_ecen662_em}:
% \begin{equation}
%   \bar S_s = (1 - \eta_s)\,\bar S_{s-1}
%   + \eta_s\,\E_{f(z \mid x^{(s)},\,\theta_{s-1})}\!\left[s\!\left(x^{(s)}, z\right)\right],
%   \label{eq:online}
% \end{equation}
% with step size $\eta_s = (s + \tau)^{-\kappa}$, $\kappa \in (\tfrac{1}{2}, 1]$.
% The M-step computes $\theta_s$ from $\bar S_s$ in closed form after each sequence.

\pagebreak

% ============================================================
\section{Experiments}
% ============================================================

We generate $S = 50$ sequences of length $T = 100$ from a $K = 3$ state HMM
with known ground truth parameters:
\begin{equation}
  \pi^* = (0.4,\; 0.35,\; 0.25), \quad
  A^* = \begin{pmatrix} 0.7 & 0.2 & 0.1 \\ 0.1 & 0.6 & 0.3 \\ 0.2 & 0.2 & 0.6 \end{pmatrix}, \quad
  \mu^* = (-3,\; 0,\; 3), \quad
  \sigma^* = (0.8,\; 0.8,\; 0.8).
\end{equation}
All four methods are initialized from the same random starting point.
We evaluate: (1) log-likelihood $\sum_s \log f_\theta(x^{(s)})$ vs.\ iteration,
(2) parameter RMSE $\|\hat\mu - \mu^*\|_2 / \sqrt{K}$ after optimal state
relabeling, and (3) the MCEM bias-variance tradeoff over
$M \in \{10, 25, 50, 100, 200, 500\}$.

% ============================================================
\section{Discussion}
% ============================================================

The four methods expose a fundamental tradeoff: Baum-Welch is exact but
assumes a tractable E-step; Variational EM introduces a structural bias via
the mean-field independence assumption; MCEM is asymptotically unbiased but
noisy for small $M$; and Online EM trades per-iteration accuracy for
scalability. Together, they illustrate that every variant of EM is maximizing
the same ELBO from Eq.~\eqref{eq:elbo}---the differences lie entirely in how
the intractable E-step is handled.

% ============================================================
\bibliographystyle{plainnat}
\bibliography{references}
\end{document}
